% \documentclass{article}
\documentclass[leqno]{article}
\usepackage[T2A]{fontenc}        
\usepackage[utf8]{inputenc}    
\usepackage[russian,english]{babel}
\usepackage{graphicx}
\usepackage{float}
\usepackage[hidelinks]{hyperref}
\usepackage{fancyvrb}
\usepackage{array}
\usepackage{booktabs}
\usepackage{siunitx}
\usepackage{tabularx}
\usepackage{longtable}
\usepackage{threeparttable}
\usepackage{amsmath}
\usepackage{mathtools}
\usepackage{bm}
% \usepackage{natbib}
\usepackage[style=numeric]{biblatex}

% \usepackage{unicode-math}
% \setmainfont{TeX Gyre Pagella}
% \setmathfont{TeX Gyre Pagella Math}

\addbibresource{learnlatex.bib} % file of reference info

\title{Computer Skills for Scientific Writing\\Математический набор текста}
\author{Абрамян Артём Арменович}
\date{Февраль 2026}

\begin{document}
  \maketitle
  \pagebreak

  \textbf{Цель работы:} Изучить возможности набора математического текста в LaTeX, математический режим LaTeX и то, как можно вводить и отображать формулы непосредственно в тексте, а также расширения, предоставляемые пакетом amsmath.\smallskip

  \section{Введение}
  \label{sec:intro}
  

  В данной лабораторной работе рассматривается математический режим LaTeX и то, как можно вводить и отображать формулы непосредственно в тексте, а также расширения, предоставляемые пакетом amsmath, и как изменять шрифты в математических выражениях.
  Набор сложных математических формул — одно из главных преимуществ LaTeX. Вы можете логически размечать математические выражения в так называемом «математическом режиме».
  
  \section{Math mode}
  \label{sec:mode}
  
  В математическом режиме пробелы игнорируются, и (почти всегда) применяется правильный интервал между символами.

  Существует два типа математического режима:
  • встроенный
  • отображаемый


    \begin{center}
      A sentence with inline mathematics: $y = mx + c$.
      A second sentence with inline mathematics:
      $5^{2}=3^{2}+4^{2}$.
      A second paragraph containing display math.
      \[
      y = mx + c
      \]
      See how the paragraph continues after the display.
    \end{center}

    Вы можете встретить математический ввод, «похожий на LaTeX», и в других местах, например, в системе MathJax для размещения уравнений на веб-страницах. Эти системы часто принимают небольшие вариации синтаксиса LaTeX, поскольку они фактически не используют LaTeX «за кулисами».
    Все наши примеры — это корректный LaTeX. Если вы видите что-то другое в другом контексте, это может быть потому, что в примере на самом деле не используется LaTeX.


  \section{Встроенный математический режим и математическая нотация}
  \label{sec:inline}
  
  Как вы видите выше, режим встроенных математических выражений обозначается парой символов доллара \(\$...\$\). Также можно использовать обозначение \( ... \). Простые выражения вводятся без какой-либо специальной разметки, и вы увидите, что математические выражения аккуратно расставлены и выделены курсивом.
  Режим встроенных математических выражений ограничивает вертикальный размер выражения, чтобы по возможности формула не нарушала межстрочный интервал абзаца.
  Обратите внимание, что все математические выражения должны быть помечены как математические, даже если это один символ.
  Используйте ... \$2\$ ... а не ... 2 ... иначе, например, когда вам нужно отрицательное число и нужно получить знак минус, ... \$-2\$ ...
  может использовать цифры, которые могут быть другого шрифта, чем цифры в тексте \(в зависимости от класса документа\). И наоборот, остерегайтесь математических конструкций, встречающихся в
  простом тексте, скопированном откуда-либо, например, денежных значениях, использующих \$, или именах файлов,
  использующих \_ \(которые могут быть помечены как \$ и \_ соответственно\).
  Мы можем легко добавлять верхние и нижние индексы; они обозначаются символами \textasciicircum{} и \_,
  соответственно.

  \begin{center}
    Superscripts $a^{b}$ and subscripts $a_{b}$.
  \end{center}


  Существует множество специализированных команд математического режима. Некоторые из них довольно просты,
  например, \\sin и \\log для синуса и логарифма или \\theta для греческой буквы.

  \begin{center}
  Some mathematics: $y = 2 \sin \theta^{2}$.
  \end{center}
  

  \section{Отображение математики}
  \label{sec:display}
  

  Для режима отображения математических выражений можно использовать те же команды, что и для встроенных вычислений.
  Режим отображения математических выражений по умолчанию центрирован и предназначен для больших уравнений,
  являющихся «частью абзаца». Обратите внимание, что среды отображения математических выражений не позволяют абзацу заканчиваться внутри математических выражений, поэтому у вас не должно быть пустых строк в исходном коде отображения.
  Абзац всегда должен начинаться перед отображением, поэтому не оставляйте пустую строку перед средой отображения математических выражений. Если вам нужно несколько строк математических выражений,
  не используйте последовательные среды отображения математических выражений (это приводит к непоследовательному
  интервалу); используйте одну из многострочных сред отображения, например, align из пакета
  amsmath, описанного далее.
  Это особенно полезно для интеграций, например:

  \begin{center}
    A paragraph about a larger equation
    \[
    \int_{-\infty}^{+\infty} e^{-x^2} \, dx
    \]
  \end{center}

  Обратите внимание, как здесь используется нижний/верхний индекс для установки пределов интегрирования.
  Мы добавили здесь один пробел вручную: \, создающий тонкий пробел перед
  dx. Форматирование оператора дифференциала различается: некоторые издатели используют прямую букву «d», 
  а другие — курсивную. Один из способов написать исходный код, позволяющий обрабатывать любой из вариантов, — 
  это создать команду \\diff, которую можно настраивать по мере необходимости, например.

  \smallskip

  \newcommand{\diff}{\mathop{}\!d} % For italic
  % \newcommand{\diff}{\mathop{}\!\mathrm{d}} % For upright
  \begin{center}
  A paragraph about a larger equation
  \[
  \int_{-\infty}^{+\infty} e^{-x^2} \diff x
  \]
  \end{center}


  Часто возникает необходимость в пронумерованном уравнении, которое создается с помощью среды для создания уравнений. Давайте повторим тот же пример:


  \begin{center}
    A paragraph about a larger equation
    \begin{equation}
    \int_{-\infty}^{+\infty} e^{-x^2} \, dx
    \end{equation}
  \end{center}

  Номер уравнения автоматически увеличивается и может быть простым числом,
  как в этом примере, или может начинаться с номера раздела, например, (2.5) для 5-го
  уравнения в разделе 2. Детали форматирования задаются классом документа
  и здесь не описываются.





  \section{Пакет amsmath}
  \label{sec:amsmath}
  
  Математическая нотация очень богата, а это значит, что инструменты, встроенные в ядро LaTeX, не могут охватить всё. 
  Пакет amsmath расширяет поддержку ядра, охватывая гораздо больше идей. 
  Руководство пользователя amsmath содержит гораздо больше примеров, чем мы можем показать в этом уроке.

  Solve the following recurrence for $ n,k\geq 0 $:
  \begin{align*}
  Q_{n,0} &= 1 \quad Q_{0,k} = [k=0]; \\
  Q_{n,k} &= Q_{n-1,k}+Q_{n-1,k-1}+\binom{n}{k},
  \quad\text{for $n$, $k>0$.}
  \end{align*}

  Окружение align* выравнивает уравнения по амперсандам и символам \&, как в таблице. 
  Обратите внимание, как мы использовали \\quad для вставки небольшого пробела и \\text для размещения обычного текста в математическом режиме. Мы также использовали другую команду математического режима, \\binom, для бинома.
  Обратите внимание, что здесь мы использовали align*, и уравнение не было пронумеровано.
  В большинстве математических окружений уравнения нумеруются по умолчанию, а вариант со звездочкой (*) отключает нумерацию.
  Пакет также имеет несколько других удобных окружений, например, для матриц.

  AMS matrices.
  \begin{center}
    \[
    \begin{matrix}
    a & b & c \\
    d & e & f
    \end{matrix}
    \quad
    \begin{pmatrix}
    a & b & c \\
    d & e & f
    \end{pmatrix}
    \quad
    \begin{bmatrix}
    a & b & c \\
    d & e & f
    \end{bmatrix}
    \]
  \end{center}

  \smallskip



  \section{Шрифты в математическом режиме}
  \label{sec:fonts}

  В отличие от обычного текста, изменение шрифта в математическом режиме часто передает очень специфический смысл.
  Поэтому они часто записываются явно. Вам понадобится следующий набор команд:
  • \\mathrm: прямой шрифт (прямой)
  • \\mathit: курсив с интервалом «текст»
  • \\mathbf: жирный шрифт
  • \\mathsf: без засечек
  • \\mathtt: моноширинный шрифт (печатная машинка)
  • \\mathbb: двойной штрих (жирный шрифт, как на доске) (предоставляется пакетом amsfonts)
  Каждая из этих команд принимает в качестве аргумента латинские буквы, поэтому, например, мы можем написать
  матрицу как

  \begin{center}
  The matrix $\mathbf{M}$.
  \end{center}

  Обратите внимание, что курсивный шрифт по умолчанию разделяет буквы, чтобы их можно было использовать для обозначения
  произведения переменных. Используйте команду \\mathit, чтобы сделать слово курсивным.
  Команды шрифта \\math.. используют шрифты, предназначенные для математических операций. Иногда вам
  необходимо встроить слово, которое является частью внешней структуры предложения и требует использования
  текущего шрифта текста; для этого вы можете использовать команду \\text{...} (которая предоставляется пакетом
  amsmath) или специальные стили шрифтов, такие как \\textrm{..}.

  \begin{center}
  $\text{bad use } size \neq \mathit{size} \neq \mathrm{size} $
  \textit{$\text{bad use } size \neq \mathit{size} \neq
  \mathrm{size} $}
  \end{center}

  \section{Дальнейшие выравнивания amsmath}
  \label{sec:alignment}

  В дополнение к среде align*, показанной в основном уроке, amsmath имеет
  несколько других конструкций для отображения математических формул, в частности gather для многострочных отображений,
  которые не требуют выравнивания, и multline для разделения большого единого выражения на
  несколько строк, выравнивая первую строку слева, а последнюю справа. Во всех случаях
  форма * по умолчанию опускает номера уравнений.

  \begin{center}
    Gather
    \begin{gather}
    P(x)=ax^{5}+bx^{4}+cx^{3}+dx^{2}+ex +f\\
    x^2+x=10
    \end{gather}
    Multline
    \begin{multline*}
    (a+b+c+d)x^{5}+(b+c+d+e)x^{4} \\
    +(c+d+e+f)x^{3}+(d+e+f+a)x^{2}+(e+f+a+b)x\\
    + (f+a+b+c)
    \end{multline*}
  \end{center}


  \section{Выравнивание столбцов в математических выражениях}
  \label{sec:colms}

  Среды выравнивания amsmath предназначены для работы с парами столбцов, где
  первый столбец каждой пары выровнен по правому краю, а второй — по левому.
  Это позволяет отображать несколько уравнений, каждое из которых выровнено по отношению к своему символу связи.

  \begin{center}
    Aligned equations
    \begin{align*}
    a &= b+1 & c &= d+2 & e &= f+3 \\
    r &= s^{2} & t &=u^{3} & v &= w^{4}
    \end{align*}
  \end{center}

  Кроме того, существуют варианты отображаемых сред, оканчивающихся на "ed", которые образуют
  подтермин внутри более крупного отображения. Например, aligned и gathered являются вариантами
  соответственно align и gather.

  \begin{center}
    \begin{itemize}
    \item
    $\begin{aligned}[t]
    a&=b\\
    c&=d
    \end{aligned}$
    \item
    $\begin{aligned}
    a&=b\\
    c&=d
    \end{aligned}$
    \end{itemize}
  \end{center}


  \section{Bold Math}
  \label{sec:bold}

  В стандартном LaTeX есть два способа выделить символы в математических выражениях жирным шрифтом. 
  Чтобы выделить жирным шрифтом всё выражение целиком, используйте команду `\\boldmath` перед вводом выражения. 
  Также доступна команда `\\mathbf` для выделения отдельных букв или слов прямым жирным шрифтом.

  \begin{center}
    $(x+y)(x-y)=x^{2}-y^{2}$
    {\boldmath $(x+y)(x-y)=x^{2}-y^{2}$ $\pi r^2$}
    $(x+\mathbf{y})(x-\mathbf{y})=x^{2}-{\mathbf{y}}^{2}$
    $\mathbf{\pi} r^2$ % bad use of \mathbf
  \end{center}

  \smallskip


  Если вам нужно получить доступ к символам, выделенным жирным шрифтом (как это делается с помощью команды \\boldmath), 
  внутри выражения, содержащего обычный вес символов, вы можете использовать команду \\bm из пакета bm. Обратите внимание, 
  что команда \\bm также работает с такими символами, как знак равенства (=) и греческие буквы.
  (Обратите внимание, что команда \\mathbf не влияет на символ \\pi в приведенном выше примере.)


  \begin{center}
    $(x+\mathbf{y})(x-\mathbf{y})=x^{2}-{\mathbf{y}}^{2}$
    $(x+\bm{y})(x-\bm{y}) \bm{=} x^{2}-{\bm{y}}^{2}$
    $\alpha + \bm{\alpha} < \beta + \bm{\beta}$
  \end{center}




  \section{Математические инструменты}
  \label{sec:tools}

  Пакет mathtools загружает amsmath и добавляет несколько дополнительных функций, таких как:
  варианты матричных окружений amsmath, позволяющие указывать выравнивание столбцов.

  \begin{center}
    \[
    \begin{pmatrix*}[r]
    10&11\\
    1&2\\
    -5&-6
    \end{pmatrix*}
    \]
  \end{center}


  \smallskip



  \section{Unicode Math}
  \label{sec:unicode}

  Существуют варианты TeX-движков, использующих шрифты OpenType. По умолчанию эти движки
  по-прежнему используют классические математические шрифты TeX, но вы можете использовать пакет unicode-math 
  для использования
  математических шрифтов OpenType. Подробности об этом пакете выходят за рамки данного курса, и мы
  отсылаем вас к документации пакета. Однако здесь мы приводим небольшой пример.

  % \begin{document}
  %   One two three
  %   \[
  %   \log \alpha + \log \beta = \log(\alpha\beta)
  %   \]
  %   Unicode Math Alphanumerics
  %   \[A + \symfrak{A}+\symbf{A}+ \symcal{A} + \symscr{A}+
  %   \symbb{A}\]
  % \end{document}


  \smallskip


  \section{Упражнения}
  \label{sec:ex}

  \subsection{Попробуйте выполнить несколько простых заданий в режиме работы с математическими выражениями: возьмите примеры и переключайтесь между встроенным и отображаемым режимами работы с математическими выражениями. Видите, какой эффект это дает?}
  \label{subsec:first}



  \begin{center}
    A sentence with inline mathematics: $y = mx + c$.

    A second paragraph containing display math.
    \[
    y = mx + c
    \]
See how the paragraph continues after the display.
  \end{center}



  \subsection{Попробуйте добавить другие греческие буквы, как строчные, так и заглавные. Вы должны суметь назвать их названия.}
  \label{subsec:second}

  \begin{center}
    $\mathbf{\pi}$
    $\Pi$
    $\theta$
  \end{center}


  \subsection{Поэкспериментируйте с командами изменения шрифта: что произойдет, если вы попытаетесь вложить их друг в друга?}
  \label{subsec:font}

  \begin{center}
    The matrix $\mathit{M}$.
    The matrix $\mathbf{\mathit{M}}$.
    The matrix $\mathbf{M}$.
  \end{center}

  \subsection{Математические формулы по умолчанию центрированы; попробуйте добавить параметр класса документа
[fleqn] (выровнять уравнение по левому краю) к некоторым из приведенных выше примеров, 
чтобы увидеть другое расположение. Аналогично, номера уравнений обычно располагаются справа.
 Поэкспериментируйте с добавлением параметра класса документа [leqno] (номера уравнений слева).}
  \label{subsec:experm}

  \begin{center}
    The matrix $\mathit{M}$.
    The matrix $\mathbf{\mathit{M}}$.
    The matrix $\mathbf{M}$.
  \end{center}
  
  \section{Заключение}
  \label{sec:conclusion}
  
  \textbf{Выводы:} В данной работе мы изучили возможности набора математического текста в LaTeX, математический режим LaTeX и то, как можно вводить и отображать формулы непосредственно в тексте, а также расширения, предоставляемые пакетом amsmath.

\end{document}